\section{Experimental Methodology}\label{sec:method}

\subsection{Evaluation Environment and Replay Dataset}
To rigorously evaluate coordination and specialization without confounding stochastic variance from external API calls, all experiments are conducted via deterministic replay through the real orchestrator hub with \texttt{EVALUATOR\_MOCK\_MODE=1}.

\noindent\textbf{Replay Stimuli and Dataset Structure.} All experiments operate on a fixed set of 64 technical-interview replay instances. Each instance provides a deterministic sequence of evaluator inputs, enabling matched comparisons across the split and unified architectures and their ablations. The 64 instances cover 8 technical interview question domains (arrays, dynamic programming, trees/graphs, system design, concurrency, databases, networking, and object-oriented design), with 8 candidate answers per domain evaluated against rubric keywords, fine-tuned cross-encoders, and concept extractors. In replay mode, the hub consumes these fixed evaluator inputs directly, driving the full orchestrator state machine and agent coordinator. Follow-up probe questions receive the candidate payload corresponding to the root concept.

\noindent\textbf{The 8-Cluster Structure.} A critical methodological detail is that replay sessions are generated by combining the 8 underlying replay clusters with 20 seeded pseudo-random question-order permutations (seeds 42--61), yielding $N = 160$ sessions per experimental condition. These 160 sessions are \textbf{not independent samples}; they reuse the identical 8 replay clusters. Presenting order permutations as independent evaluation sessions would artificially deflate variance and produce overconfident inferences. We therefore explicitly model the 8 underlying replay clusters as the independent structural units in all statistical analyses.

\noindent\textbf{Evaluation Mode.} All experimental evidence is generated via deterministic replay through the real orchestrator hub, testing coordination and trajectory dynamics under fixed stimuli. No live generative LLM is invoked; the LLM gateway is stubbed as down, and deterministic template renderers handle text output, ensuring that trajectory differences are attributable solely to the evaluated coordination and agent architectures.

\subsection{Baseline and Ablation Matrix}
To address RQ1 (coordination dynamics) and RQ3 (coordination mechanisms), we evaluate the baseline platform alongside controlled ablation conditions across the identical 160 replay sessions:
\begin{itemize}
\item \textbf{Frozen Baseline ($B_0$):} The original single-controller workflow without agents. Difficulty transitions are chosen solely by the frozen PPO policy (or rule fallback) bounded by guardrails; follow-up probes are governed by the static hub threshold.
\item \textbf{Full Multi-Agent System ($M_1$):} The complete coordinated architecture with Strategy Agent, Coaching Agent, peer messaging, authority enforcement, and threshold-gated arbitration ($\tau = 0.6$).
\item \textbf{Silent Agent ($M_2$):} The multi-agent layer is active, but agents emit zero proposals, testing whether the coordination infrastructure is completely transparent when silent.
\item \textbf{Coaching Ablated ($M_3$):} The Coaching Agent is disabled while Strategy remains active, testing whether difficulty trajectories are independent of coaching state.
\item \textbf{Arbitration Ablated ($M_4$):} Arbitration checks (confidence threshold, evidence requirement, one-step bound) are removed, admitting raw Strategy proposals.
\item \textbf{Peer Messaging Ablated ($M_5$):} Strategy proposals are not delivered to Coaching as peer messages, isolating the effect of inter-agent information exchange.
\item \textbf{PPO Policy Ablated ($M_7$):} The learned PPO policy is disabled, falling back to a deterministic heuristic rule to evaluate interaction between Strategy heuristics and the baseline.
\item \textbf{Canonical Guardrails ($B_{0c}$):} Baseline workflow evaluated with canonical \texttt{rl/guardrails.py} instead of deployed inline rules.
\end{itemize}

\subsection{Framework for Evaluating Agent Specialization (RQ2)}
To evaluate whether decomposing technical coaching into specialized agents provides measurable value over a unified controller (RQ2), we formulate a controlled comparative framework contrasting:
\begin{enumerate}
\item \textbf{Specialized Architecture ($M_{\text{split}}$):} The two-agent baseline ($M_1$) partitioning responsibilities across Strategy and Coaching with disjoint state ledgers and explicit peer messaging.
\item \textbf{Unified Architecture ($M_{\text{unified}}$):} A monolithic agent (\texttt{UnifiedCoachAgent}) maintaining a joint session ledger and synthesizing difficulty actions, follow-up probes, and coaching plans in a single dispatch.
\end{enumerate}

\noindent\textbf{Equivalence Conditions.} To ensure a fair comparison, both configurations receive identical observation inputs, identical question sequences, identical heuristic decision thresholds, and identical arbitration rules. We evaluate five operational outcomes: (1)~decision agreement rate ($\mathbb{P}(\text{Decision}_{\text{unified}} = \text{Decision}_{\text{split}})$); (2)~difficulty trajectory divergence; (3)~follow-up probe selection; (4)~per-turn dispatch and serialization latency; and (5)~architectural state encapsulation (schema separation and ledger isolation).

\subsection{Threshold Decomposition (Experiment XC)}
Because the unarbitrated condition ($M_4$) simultaneously removes the confidence threshold, evidence verification, and one-step limits, it cannot isolate which mechanism prevents oscillation collapse. We therefore execute Experiment XC, sweeping the confidence threshold $\tau \in \{0.0, 0.5, 0.6, 0.7, 0.8, 0.9\}$ across the same 160 replay sessions while holding all other safety envelopes and evidence checks active.

\subsection{Statistical Methodology}
For session-level point estimates and paired differences, we compute standard non-parametric percentile bootstrap intervals with $B = 10{,}000$ resamples. To account for the 8-cluster dependency, we implement a cluster bootstrap over the 8 underlying replay clusters following Field and Welsh~\cite{fw07} ($B = 10{,}000$), resampling entire replay clusters with replacement. In addition, we report a leave-one-cluster-out (LOCO) sensitivity analysis to verify whether observed effects depend on individual influential replay clusters. All reported bootstrap confidence intervals are descriptive measures of parameter dispersion across resamples; no formal null-hypothesis significance tests or $p$-values are evaluated.
